\chapter{Visual Perception and Deep Learning}
\label{ch:vision-deep-learning}

\textit{Visual perception feels effortless from the inside, but the computational problem is severe. The brain must start from a rapidly changing pattern of pixel-like sensory measurements and recover stable facts about objects, places, motion, and action. This lecture frames visual recognition as a problem of finding the right representation, then uses deep learning to show how such representations can be learned and compared to the ventral visual stream.}

\section{Why vision is a higher-brain problem}
\label{sec:vision-higher-brain}

The lecture begins with a simple question: what does vision need to recover from the retinal image? At minimum, the system wants to know what is out there, where it is, and what it is likely to do next. Those goals already point beyond raw sensation. A camera records local luminance values, but perception requires interpretation: grouping edges into objects, deciding whether two views belong to the same thing, and turning ambiguous input into a coherent guess about the world.

This is why the course treats perception as part of higher brain function rather than as a passive readout of the environment. The visual system must infer latent causes from incomplete sensory evidence. It is not enough to detect brightness changes; the system must construct internal variables that stay useful across changes in viewpoint, position, scale, illumination, and background clutter.

\begin{aside}[Moravec's paradox]
The lecture invokes Moravec's paradox to make a useful point: tasks that feel easy to humans, such as recognizing a familiar object across wildly different viewing conditions, are often the tasks that proved hardest to automate. By contrast, tasks that feel formally difficult, such as symbolic calculation, can be much easier for machines once the rules are explicit.
\end{aside}

\section{Invariant object recognition is the core challenge}
\label{sec:invariant-recognition}

Object recognition is hard because the same object can generate many different retinal images. Translating an object across the visual field changes many pixel values at once. Rotation, distance, occlusion, lighting, and background context change them again. In raw pixel space, images of the same category do not sit in one tidy cluster. Instead, each object class traces out a high-dimensional manifold of possible appearances, and those manifolds can be heavily interleaved \parencite{dicarlo2007untangling}.

The key phrase from the lecture is \emph{invariant object recognition}. A good visual representation should preserve object identity while being tolerant to nuisance variables. That does not mean ignoring all variation. Position, pose, and illumination still matter for behavior. It means that small natural changes should produce smooth changes in representation, while category identity should become easier to read out.

Another way to say this is that visual cortex must perform a disentangling computation. Early sensory signals are in a bad ``pixel'' format for recognition. Recognition becomes easier only after the system transforms those signals into features in which relevant variables are more explicit and category boundaries are less tangled.

\begin{figure}[t]
  \centering
  \begin{tikzpicture}[>=Latex, font=\small]
    \draw[rounded corners=4pt, thick] (0,0) rectangle (4.8,3.4);
    \draw[rounded corners=4pt, thick] (7.2,0) rectangle (12,3.4);
    \node[font=\sffamily\bfseries] at (2.4,3.75) {Pixel-like space};
    \node[font=\sffamily\bfseries] at (9.6,3.75) {Learned feature space};

    \draw[thick, black!60] (1.6,1.8) ellipse (1.3 and 0.85);
    \draw[thick, black!60] (2.8,1.45) ellipse (1.3 and 0.85);
    \node at (0.9,2.7) {object A};
    \node at (3.7,0.75) {object B};

    \draw[thick, black!60] (8.6,2.1) ellipse (1.0 and 0.7);
    \draw[thick, black!60] (10.8,1.0) ellipse (1.0 and 0.7);
    \node at (8.5,3.0) {object A};
    \node at (10.8,2.0) {object B};

    \draw[->, thick] (5.3,1.7) -- (6.7,1.7);
    \node[align=center] at (6,2.4) {hierarchical\\representation\\learning};
  \end{tikzpicture}
  \caption{The lecture's basic geometric intuition. In raw sensory coordinates, category manifolds are tangled. A useful visual representation transforms them into a space where identity is easier to read out with a simple downstream decision.}
\end{figure}

The figure captures the main conceptual move of the lecture: deep networks matter not because they magically solve vision, but because they provide a mechanism for gradually reshaping the representation until simple decisions become possible.

\section{Supervised deep learning made this concrete}
\label{sec:supervised-learning}

The lecture uses the ImageNet era as the turning point. ImageNet provided a very large supervised benchmark: on the order of 15 million labeled images spanning 1000 categories. That benchmark forced models to solve recognition at realistic visual scale rather than on toy datasets. AlexNet's 2012 win demonstrated that deep convolutional networks could dramatically outperform the previous field on this task, reducing the top-5 validation error to 16.4\% while the next-best system remained above 26\% \parencite{krizhevsky2012imagenet}.

That result was important for two reasons. First, it showed that the right combination of depth, nonlinearities, optimization, and data could learn useful hierarchical visual features rather than relying on hand-engineered descriptors. Second, it gave neuroscience a new family of concrete models to compare against cortex. For the first time, researchers had trainable systems that solved a hard perceptual task well enough to make their internal representations interesting.

The lecture also places AlexNet in a broader architectural trajectory. Later systems such as VGG, GoogLeNet, and ResNet deepened the hierarchy and improved optimization. Residual networks are especially relevant because their skip connections make it easier to train very deep models: instead of forcing a block to learn an entirely new transformation, the block can learn a correction to the identity map \parencite{he2016deep}. This is one way to combat optimization pathologies such as vanishing gradients and depth-related degradation.

\section{What a deep network computes}
\label{sec:deep-network-computation}

At the mathematical level, a deep network is a stack of linear-nonlinear transformations. Each layer receives an input vector, applies an affine map, and passes the result through a nonlinearity:

\begin{keyequation}[Layer computation]
\begin{equation}
  \vec{y} = \phi\!\left(W \vec{x} + \vec{b}\right).
\end{equation}
\tcblower
$\vec{x}$ is the incoming activity pattern, $W$ is the weight matrix, $\vec{b}$ is a bias term, and $\phi$ is a nonlinear activation such as a rectified linear unit.
\end{keyequation}

If the network were only a stack of linear maps, the composition would still be linear and nothing genuinely new would happen with depth. The nonlinearity is what allows the network to progressively carve up input space and build increasingly abstract features. In vision, early filters may resemble edge detectors or local contrast detectors, later filters respond to textures and parts, and still later layers become useful for category-level distinctions.

\subsection{Why convolutions help}

Convolutional layers encode two priors that are especially well matched to images. First, they use \emph{local receptive fields}: a filter only looks at a small patch at a time, reflecting the fact that meaningful local structure is spatially organized. Second, they use \emph{weight sharing}: the same filter is applied at many image locations, which means the model can detect the same pattern wherever it appears. Pooling and striding then trade fine spatial detail for a larger effective receptive field and more tolerance to small translations.

These choices do not make a network fully invariant. Instead, they bias the model toward learning representations where small shifts and deformations matter less than stable feature combinations. That is exactly the representational move the lecture motivates from the object-recognition problem.

\subsection{How learning works}

The lecture also reviews backpropagation as the mechanism that makes deep supervised learning possible. The network is evaluated with a loss function that compares its output to the desired label, and the chain rule is used to compute how each parameter contributed to the error. Parameters are then nudged to reduce the loss:

\begin{keyequation}[Gradient-descent update]
\begin{equation}
  w \leftarrow w - \eta \frac{\partial \mathcal{L}}{\partial w}.
\end{equation}
\tcblower
$\mathcal{L}$ is the loss, $w$ is a parameter, and $\eta$ is the learning rate. Backpropagation efficiently computes the needed gradients across many layers.
\end{keyequation}

Backpropagation gives the network a way to assign credit across depth, but it also highlights one of the lecture's recurring caveats. The algorithm is extremely successful as an engineering procedure, yet several of its ingredients are biologically questionable: exact error backpropagation, massive labeled datasets, and the degree of weight sharing used by convolutions are not straightforward mirrors of cortical learning.

\section{Deep networks as models of the ventral stream}
\label{sec:deep-models-ventral-stream}

Once deep networks became good at categorization, the natural next question was whether their internal representations resembled those in visual cortex. The lecture emphasizes that the comparison cannot be done neuron by neuron. A hidden unit in a convolutional net is not uniquely aligned to a particular biological neuron, and even repeated trainings of the same architecture need not preserve a one-to-one correspondence between units.

Representational similarity analysis (RSA) solves this by comparing \emph{geometry} rather than individual units \parencite{kriegeskorte2008representational}. For a set of stimuli, one computes the response vector to each stimulus and then asks how dissimilar each pair of responses is. The result is a representational dissimilarity matrix:

\begin{keyequation}[Representational dissimilarity matrix]
\begin{equation}
  \mathrm{RDM}_{ij} = 1 - \mathrm{corr}(r_i, r_j).
\end{equation}
\tcblower
Here $r_i$ and $r_j$ denote the activity patterns evoked by stimuli $i$ and $j$. Similar representations give similar RDM structure, even if the underlying units are not directly matched.
\end{keyequation}

This framework made it possible to compare fMRI data from human inferotemporal cortex, electrophysiological recordings from monkey IT, and hidden-layer representations from deep networks on the same footing. The central finding highlighted in the lecture is that supervised deep models capture the geometry of IT activity substantially better than unsupervised alternatives \parencite{khalighrazavi2014deep}. Goal-driven deep models that categorize well also tend to predict neural response structure better \parencite{yamins2016using}.

\begin{aside}[What counts as ``similar''?]
The claim is not that the brain literally contains AlexNet. The claim is narrower and more useful: when a model is trained to solve the same task the brain solves, some of the representational geometry it discovers resembles the geometry observed in ventral-stream data.
\end{aside}

\section{Deep networks as experimental tools}
\label{sec:deep-network-tools}

The lecture goes one step further than comparison and uses deep networks as tools for experimental control. The Bashivan et al. example proceeds in four stages \parencite{bashivan2019neural}:
\begin{enumerate}[leftmargin=*]
  \item train a convolutional network on a large image-recognition task,
  \item fit a differentiable mapping from network features to recorded V4 responses,
  \item optimize the input image to maximize the activity of a chosen neuron or neural population,
  \item present the synthesized image to the animal and measure whether the predicted control is achieved.
\end{enumerate}

This is conceptually important. If a model has learned something close to a neuron's preferred feature geometry, then optimizing through the model should generate images that strongly drive that neuron. The lecture shows exactly that: CNN-optimized images can strongly activate V4 sites, and more sophisticated optimization can even approximate one-hot population control. In other words, the model is not just a classifier; it becomes a hypothesis generator about what the neural code is sensitive to.

That move is especially attractive in systems neuroscience because it links representation learning to intervention. Instead of only asking whether a network looks similar to cortex, we can ask whether it helps us design stimuli that probe cortical computations more precisely.

\section{Failures are scientifically informative}
\label{sec:failures-limits}

The lecture does not present deep learning as a finished theory of vision. It closes by listing several reasons for caution. Architecture design is still partly trial-and-error. Deep networks can be fooled by adversarial perturbations: tiny, often human-imperceptible changes to an image can force an arbitrary classification \parencite{goodfellow2015explaining}. That fragility shows that good benchmark performance does not automatically imply human-like visual understanding.

There are also open biological questions. Standard deep networks are usually feedforward at inference time, while cortex is richly recurrent. Training relies on large supervised datasets, while biological organisms learn from far less explicit labeling. Convolution imposes strict translational weight sharing, whereas cortical circuitry is only approximately repetitive. These mismatches do not make deep networks useless for neuroscience, but they do limit how literally we should interpret them as brain models.

The productive reading of the lecture is therefore balanced. Deep networks are powerful because they turn representational hypotheses into trainable systems. They are compelling because they solve real visual tasks and partially align with cortical data. But they are still models: abstractions that illuminate some aspects of vision while distorting others.

\section{Take-home view}
\label{sec:take-home-view}

The lecture's central claim can be summarized in one sentence: visual recognition becomes tractable when the system learns a representation in which object identity is easier to separate from nuisance variation. Deep learning matters because it provides a concrete recipe for constructing such representations, evaluating them on difficult tasks, and comparing them to cortical population codes.

That is why deep nets occupy a double role in the course. They are engineering tools for object recognition, and they are scientific instruments for thinking about higher brain function. The same network can be judged by three standards at once: how well it categorizes images, how brain-like its internal geometry is, and how useful it is for generating new experimental questions.

\begin{chaptersummary}

\begin{center}
\renewcommand{\arraystretch}{1.6}
\begin{tabularx}{\textwidth}{@{}l X X@{}}
  \toprule
  \textbf{Concept} & \textbf{Equation or idea} & \textbf{Key insight} \\
  \midrule
  Invariant recognition &
  Same object, many retinal images &
  Recognition requires tolerance to nuisance variation \\
  %
  Disentangling &
  Transform tangled manifolds into separable features &
  Good representations make simple readout possible \\
  %
  Deep layer &
  $\vec{y} = \phi(W \vec{x} + \vec{b})$ &
  Depth plus nonlinearity builds increasingly useful features \\
  %
  Backpropagation &
  $w \leftarrow w - \eta \, \partial \mathcal{L} / \partial w$ &
  Learning assigns credit across many layers \\
  %
  RSA &
  $\mathrm{RDM}_{ij} = 1 - \mathrm{corr}(r_i, r_j)$ &
  Compare representational geometry without matching units one by one \\
  %
  Adversarial examples &
  Tiny perturbations, large classification changes &
  High performance does not guarantee human-like robustness \\
  \bottomrule
\end{tabularx}
\end{center}

\end{chaptersummary}

\begin{conceptquestions}
  \item Why is raw pixel space a poor format for object recognition? Give two examples of natural changes that preserve object identity while strongly changing pixel values.

  \item The lecture says that a good visual representation should change smoothly with position, illumination, or angle. Why is that property useful for recognition?

  \item AlexNet's ImageNet result was treated as a turning point. What did it show about the role of depth, data, and learned features in visual recognition?

  \item Convolutional layers use local receptive fields and weight sharing. How do those two design choices help with natural images, and what limitations do they leave unresolved?

  \item Why is nonlinearity necessary in a deep network? What would be lost if every layer were purely linear?

  \item Backpropagation is an effective learning rule for artificial networks. Why does the lecture still describe it as biologically questionable?

  \item What problem does representational similarity analysis solve when comparing deep networks to cortical data? Why is a unit-by-unit comparison usually not meaningful?

  \item Supervised models often match IT representations better than unsupervised models. What does that suggest about the relationship between task demands and learned representations?

  \item In the V4 optimal-stimulus experiments, why is it useful that the mapping from image to predicted neural response is differentiable?

  \item The lecture ends with both enthusiasm and caution about deep learning. What are the strongest reasons to treat deep nets as useful scientific models, and what are the strongest reasons not to mistake them for full theories of vision?
\end{conceptquestions}
